\chapter{COGARCH Model and Estimation}

\section{Model Specification}
The continuous‑time jump‑diffusion COGARCH model \cite{kluppelberg2004continuous} is defined by
\begin{align}
\frac{dS_t}{S_{t-}} &= \alpha\,dt + \sqrt{V_{t-}}\,dW_t + \bigl(e^{\sqrt{V_{t-}}Z} - 1\bigr)\,dN_t , \label{eq:price_sde}\\
dV_t &= \kappa(\theta - V_{t-})\,dt + \phi\, V_{t-}\,Z^2\,dN_t , \label{eq:var_sde}
\end{align}
where $W_t$ is a standard Brownian motion, $N_t$ a Poisson process with intensity $\lambda$, and $Z\sim\mathcal{N}(0,1)$ i.i.d.\ jump sizes. The key parameter for criticality is the branching ratio
\begin{equation}
B = \frac{\lambda\,\phi\,\mathbb{E}[Z^2]}{\kappa}. \label{eq:B}
\end{equation}
When $B \lesssim 1$, the volatility process exhibits near‑critical behaviour: long memory, heavy tails, and multifractal scaling.

\section{Method of Simulated Moments (MSM)}
Because the likelihood of the COGARCH model is intractable, we employ the Method of Simulated Moments. The following moments are matched:
\begin{itemize}
    \item Mean, standard deviation, skewness, excess kurtosis of returns;
    \item Autocorrelation of $|r_t|$ at lags 1 and 5;
    \item Hill tail index (top 5\% of $|r_t|$).
\end{itemize}
The objective function minimises the weighted squared distance between empirical moments and their simulated counterparts. We use differential evolution with 10 simulations per evaluation to avoid local optima.

\section{Estimation Results}
Table \ref{tab:params} reports the estimated daily discrete‑time parameters (derived from the exact discretisation of the COGARCH) and the annualised continuous‑time parameters.

\begin{table}[H]
\centering
\caption{Estimated MSM parameters for the COGARCH model (post‑break sample).}
\label{tab:params}
\begin{tabular}{l c}
\toprule
\textbf{Parameter} & \textbf{Estimate} \\
\midrule
$\kappa$ (mean reversion speed, annual) & $0.0729$ \\
$\theta$ (long‑run variance, annual) & $0.000213$ \\
$\phi$ (volatility jump sensitivity) & $0.7100$ \\
$\lambda$ (daily jump probability) & $0.0767$ \\
$\lambda$ (annual jump intensity) & $19.335$ \\
$\mathbb{E}[Z^2]$ & $1.000$ \\
\bottomrule
\end{tabular}
\end{table}

The branching ratio computed from \eqref{eq:B} using the daily $\lambda$ (since the discrete approximation uses daily probability) is
\[
B = \frac{0.0767 \times 0.7100 \times 1}{0.0729} = 0.7468.
\]
This value is below the near‑critical window $[0.9,1)$ typically observed in developed markets \cite{hardiman2013critical}. However, as we show next, the model still generates long memory and multifractality.


